% !TeX root = main.tex
\changecolours{ScoutNavy}
\section[Modelagem: Embedding vs. Referencing]{Teoria de Schema Design: \textit{Embedding}, \textit{Referencing} e Padrões}

% ------------------------------------------------------------------------------
\twocolframe[title={Princípios de Schema Design no MongoDB},leftcol=7.0cm,rightcol=7.0cm]{%
  \alert{Três Estratégias Fundamentais}\par
  \itemseps[topsepi=0.2mm,itemsepi=0.5mm,labelsep=1.5mm,leftmargini=2.5mm]
  \scriptsize
  \begin{enumerate}
    \item \textbf{\textit{Embedded Data} (Dados Incorporados):} Aninhar objetos relacionados dentro do mesmo documento (ex.: um livro com o subdocumento \texttt{author}, ou uma cliente com a lista de \texttt{enderecos}).
    \item \textbf{\textit{References} (Referências):} Guardar o \texttt{ObjectId} de um documento em outra coleção (ex.: coleção \texttt{courses} apontando para \texttt{professor: ObjectId(...)}).
    \item \textbf{\textit{Denormalization} (Desnormalização):} Copiar campos específicos (como \texttt{authorName}) para evitar \texttt{\$lookup} em leituras críticas.
  \end{enumerate}
}{%
  \alert{A Métrica-Chave: \textit{Read-Write Ratio}}\par
  \itemseps[topsepi=0.2mm,itemsepi=0.5mm,labelsep=1.5mm,leftmargini=2.5mm]
  \scriptsize
  \begin{itemize}
    \item \textbf{Proporção Leitura vs. Escrita (\textit{Read-Write Ratio}):}
      \begin{itemize}\scriptsize
        \item \textbf{Cargas Intensivas em Leitura (\textit{Read-Heavy}):} Catálogos de produtos e vitrines têm muito mais leituras que escritas $\to$ favorecem \textbf{\textit{Embedding}} e \textbf{Desnormalização} (1 única busca traz tudo).
        \item \textbf{Cargas Intensivas em Escrita (\textit{Write-Heavy}):} Eventos frequentes ou entidades compartilhadas que mudam o tempo todo $\to$ favorecem \textbf{\textit{Referencing}}.
      \end{itemize}
    \item \textbf{Regra de Ouro:} \textit{``Dados lidos juntos pela aplicação devem ser armazenados juntos!''}
  \end{itemize}
}

% ------------------------------------------------------------------------------
\begin{frame}{Visualizando as Duas Estratégias de Modelagem}
  \vspace{0.15cm}
  \makebox[\textwidth][c]{%
    \begin{tikzpicture}
      % Bloco 1: Embedding (Esquerda)
      \node[draw=ScoutGreen, fill=ScoutGreen!10, very thick, rounded corners=6pt, minimum width=6.8cm, minimum height=5.0cm] (emb) at (0,0) {};
      \node[font=\bfseries\sffamily\footnotesize, text=ScoutGreen] at (0, 2.1) {1. Incorporação (\textit{Embedding} / Desnormalizado)};
      \node[draw=ScoutGreen!80!black, fill=white, thick, rounded corners=4pt, text width=6.1cm, inner sep=6pt, align=left, font=\ttfamily\fontsize{6.2pt}{7.8pt}\selectfont] at (0, -0.25) {%
        \textbf{\sffamily\color{ScoutGreen}Coleção: clientes (1 Único Documento)}\\[2pt]
        \{\\
        \hspace*{2mm}"\_id": \{ "\$oid": "671ec880...0010" \},\\
        \hspace*{2mm}"nome": "Ana Costa",\\
        \hspace*{2mm}"cpf": "123.456.789-00",\\
        \hspace*{2mm}\textbf{"enderecos": [} \textcolor{ScoutGreen}{// Subdocumentos!}\\
        \hspace*{4mm}\{ "apelido": "Casa", "rua": "Av. Chermont" \},\\
        \hspace*{4mm}\{ "apelido": "IFCE", "rua": "CE-187, Taua" \}\\
        \hspace*{2mm}\textbf{]}\\
        \}%
      };

      % Bloco 2: Referencing (Direita)
      \node[draw=ScoutNavy, fill=ScoutNavy!10, very thick, rounded corners=6pt, minimum width=6.8cm, minimum height=5.0cm] (ref) at (7.3,0) {};
      \node[font=\bfseries\sffamily\footnotesize, text=ScoutNavy] at (7.3, 2.1) {2. Referenciamento (\textit{Referencing} / Normalizado)};

      \node[draw=ScoutNavy!80!black, fill=white, thick, rounded corners=4pt, text width=6.1cm, inner sep=5.5pt, align=left, font=\ttfamily\fontsize{6.2pt}{7.6pt}\selectfont] (docA) at (7.3, 0.85) {%
        \textbf{\sffamily\color{ScoutNavy}Coleção: avaliacoes}\\[2pt]
        \{\\
        \hspace*{2mm}"\_id": \{ "\$oid": "671ec990...0999" \},\\
        \hspace*{2mm}\textbf{\color{red!80!black}"produto\_id": \{ "\$oid": "671ec880...0001" \}},\\
        \hspace*{2mm}"autor": "Lucas", "nota": 5\\
        \}%
      };

      \node[draw=ScoutNavy!80!black, fill=white, thick, rounded corners=4pt, text width=6.1cm, inner sep=5.5pt, align=left, font=\ttfamily\fontsize{6.2pt}{7.6pt}\selectfont] (docB) at (7.3, -1.55) {%
        \textbf{\sffamily\color{ScoutNavy}Coleção: produtos}\\[2pt]
        \{\\
        \hspace*{2mm}\textbf{\color{red!80!black}"\_id": \{ "\$oid": "671ec880...0001" \}},\\
        \hspace*{2mm}"sku": "NOTE-ACER-15", "preco": 4599.90\\
        \}%
      };

      \draw[->, ultra thick, red!80!black] ([xshift=-1.4cm]docA.south) -- node[right=4pt, font=\bfseries\sffamily\scriptsize, text=red!80!black] {Referência por \texttt{ObjectId}} ([xshift=-1.4cm]docB.north);
    \end{tikzpicture}%
  }
\end{frame}

% ------------------------------------------------------------------------------
\begin{frame}{O Anti-Padrão \textit{Unbounded Arrays} e os Padrões de Projeto}
  \vspace{0.05cm}
  \begin{columns}[t]
    \begin{column}{0.46\textwidth}
      \alert{Por que Evitar Arrays Ilimitados?}\par
      \itemseps[topsepi=0.2mm,itemsepi=0.4mm,labelsep=1.5mm,leftmargini=2.5mm]
      \scriptsize
      \begin{itemize}
        \item Embutir listas sem fim (ex.: 50.000 avaliações no produto) estoura o limite de \textbf{16 MB} (\texttt{BSONObjectTooLarge}) e degrada o \textit{WiredTiger}.
        \item \textbf{Cardinalidade:} \textbf{1:1 / 1:Few ($<50$)} $\to$ \textit{Embedding}; \textbf{1:Many} $\to$ \textit{Referencing} / \textit{Subset}; \textbf{1:Milhões} $\to$ Ref. no filho.
      \end{itemize}
    \end{column}
    \begin{column}{0.50\textwidth}
      \alert{Padrões de Modelagem (\textit{Design Patterns})}\par
      \vspace{0.03cm}
      \scriptsize
      \renewcommand{\arraystretch}{1.08}
      \begin{tabular}{p{1.65cm} p{2.65cm}}
        \toprule
        \textbf{Padrão} & \textbf{Aplicação Prática} \\
        \midrule
        \textbf{Polymorphic} & Categorias distintas na mesma coleção. \\
        \textbf{Subset} & Últimas avaliações no produto. \\
        \textbf{Extended Ref} & Campos do cliente/item no pedido. \\
        \textbf{Computed} & Média e total pré-calculados. \\
        \bottomrule
      \end{tabular}
    \end{column}
  \end{columns}
\end{frame}
